\chapter*{Introduction générale}
\addcontentsline{toc}{chapter}{Introduction générale}
\markboth{Introduction générale}{Introduction générale}
\label{chap:introduction}

Reconnaître une activité à partir d'un accéléromètre ou d'un gyroscope dans un
téléphone ou une montre est devenu courant, que ce soit pour le sport, la santé
ou l'aide à domicile. La manière habituelle de procéder reste simple : on
collecte des données en laboratoire, on entraîne un réseau une fois, puis on
déploie. Dès qu'un nouvel utilisateur arrive, que le capteur change de place
ou qu'une activité inédite apparaît, les performances chutent. Tout réapprendre
depuis zéro coûte cher ; un finetuning naïf fait souvent oublier ce qui
marchait. Et la plupart des systèmes se contentent de classer l'instant
présent, sans prévenir une transition qui s'annonce.

Ce mémoire part de ce constat. Nous construisons un système de HAR capable
d'apprendre au fil de l'eau quand de nouveaux utilisateurs ou de nouvelles
classes arrivent, et capable aussi de détecter tôt un changement d'activité.
L'implémentation vit dans le dépôt \texttt{har\_project} : un backbone
Transformer, un tampon de rejeu guidé par l'incertitude (UWR), une
anticipation multi-classe via une tête LSTM sur fenêtres partielles, un
adaptateur de domaine léger et un module de risque postural.

La première partie (chapitres~1 à~3) rappelle les bases de la HAR, les idées
de l'apprentissage continu et les travaux proches. La seconde (chapitres~4
à~6) décrit la conception, le code et les expériences, avec des comparaisons
à EWC, iCaRL et au finetuning séquentiel. La conclusion reprend les chiffres
et les limites.

\section{Problématique détaillée}

\subsection{Limites du paradigme batch}

Collecter un corpus, entraîner offline et figer les poids suppose une
distribution à peu près stable. En déploiement, ce n'est presque jamais le
cas. Les gens bougent différemment, le téléphone n'est pas toujours au même
endroit, de nouvelles activités apparaissent, et les signaux dérivent avec le
temps. Réentraîner périodiquement sur tout l'historique est coûteux, et les
données brutes ne peuvent pas toujours rester sur un serveur.

\subsection{Oubli catastrophique en pratique}

Quand on fine-tune sujet par sujet --- comme dans notre protocole à 44 tâches
--- l'accuracy tombe d'environ 91~\% (entraînement joint) à 62~\% sans
mécanisme anti-oubli (chapitre~6). EWC et iCaRL aident, mais ne ciblent pas
particulièrement les exemples près des frontières de décision. UWR oriente un
tampon de 2\,000 fenêtres vers celles où le modèle est le plus incertain.

\subsection{Anticipation vs classification instantanée}

Beaucoup de travaux HAR cherchent surtout une bonne accuracy sur la fenêtre
courante. Pour la sécurité (chute, sortie d'une zone sûre), il faudrait
parfois agir avant que le mouvement soit terminé. La littérature vidéo
\parencite{ryoo2011,gammulle2019} parle d'early recognition ; côté IMU, peu
d'articles mêlent anticipation et apprentissage continu. Dans
\texttt{har\_project}, on prédit la classe suivante à partir de $S=5$
fenêtres tronquées à $p\in\{25,50,75\}\%$, en ne gardant que les transitions.
On obtient un macro-F1 de validation d'environ 0,60--0,64 (baseline majorité
$\approx$0,04). Un score binaire de risque postural complète ce volet
(macro-F1 $\approx$0,88).

\section{Approche proposée}

Le modèle \texttt{HARContinualModel} s'appuie sur un encodeur Transformer
pré-entraîné (embeddings 128-D via un token CLS ; une initialisation SSL est
possible). L'apprentissage continu repose sur UWR, des prototypes de classes
et une perte contrastive, en scénarios user- et class-incremental. Autour de
ce noyau, on trouve une tête d'anticipation LSTM, un \texttt{DomainAdapter},
le fall-risk, une calibration de confiance en ligne et un petit bandit de
seuil. Chaque brique peut être testée seule ; les ablations du chapitre~6
s'en servent.

\section{Organisation du document}

Le chapitre~1 introduit le vocabulaire HAR et les pipelines classiques. Le
chapitre~2 formalise l'apprentissage continu. Le chapitre~3 compare les
travaux récents et pointe ce qui manque encore. Les chapitres~4 à~6
couvrent la conception, l'implémentation Python et l'évaluation. Les annexes
rassemblent datasets, hyperparamètres, extraits de code et fiches de lecture.
